\documentclass[aspectratio=169,11pt]{beamer}
\newcommand{\footerlabel}{DEMO101 Algorithms \textbullet{} Lecture 1}
% 共用样式 —— 虚构示例课件（DEMO），内容全部自编，不来自任何真实课程。
% 用法：每份 deck 先 \documentclass，再定义 \footerlabel，最后 % 共用样式 —— 虚构示例课件（DEMO），内容全部自编，不来自任何真实课程。
% 用法：每份 deck 先 \documentclass，再定义 \footerlabel，最后 % 共用样式 —— 虚构示例课件（DEMO），内容全部自编，不来自任何真实课程。
% 用法：每份 deck 先 \documentclass，再定义 \footerlabel，最后 \input{_style}。
\usepackage[T1]{fontenc}
\usepackage{lmodern}
\usepackage{amsmath,amssymb}
\usetheme{default}
\setbeamertemplate{navigation symbols}{}
\setbeamercolor{frametitle}{fg=blue!50!black}
\setbeamerfont{frametitle}{series=\bfseries,size=\Large}
\setbeamercolor{title}{fg=blue!50!black}
\setbeamercolor{footline}{fg=black!60}
\setbeamertemplate{footline}{%
  \leavevmode\hbox{%
    \begin{beamercolorbox}[wd=.85\paperwidth,ht=2.5ex,dp=1ex,left,leftskip=1em]{footline}\footnotesize\footerlabel\end{beamercolorbox}%
    \begin{beamercolorbox}[wd=.15\paperwidth,ht=2.5ex,dp=1ex,right,rightskip=1em]{footline}\footnotesize\insertframenumber\end{beamercolorbox}}%
  \vskip0pt}
\setbeamertemplate{itemize items}[circle]
。
\usepackage[T1]{fontenc}
\usepackage{lmodern}
\usepackage{amsmath,amssymb}
\usetheme{default}
\setbeamertemplate{navigation symbols}{}
\setbeamercolor{frametitle}{fg=blue!50!black}
\setbeamerfont{frametitle}{series=\bfseries,size=\Large}
\setbeamercolor{title}{fg=blue!50!black}
\setbeamercolor{footline}{fg=black!60}
\setbeamertemplate{footline}{%
  \leavevmode\hbox{%
    \begin{beamercolorbox}[wd=.85\paperwidth,ht=2.5ex,dp=1ex,left,leftskip=1em]{footline}\footnotesize\footerlabel\end{beamercolorbox}%
    \begin{beamercolorbox}[wd=.15\paperwidth,ht=2.5ex,dp=1ex,right,rightskip=1em]{footline}\footnotesize\insertframenumber\end{beamercolorbox}}%
  \vskip0pt}
\setbeamertemplate{itemize items}[circle]
。
\usepackage[T1]{fontenc}
\usepackage{lmodern}
\usepackage{amsmath,amssymb}
\usetheme{default}
\setbeamertemplate{navigation symbols}{}
\setbeamercolor{frametitle}{fg=blue!50!black}
\setbeamerfont{frametitle}{series=\bfseries,size=\Large}
\setbeamercolor{title}{fg=blue!50!black}
\setbeamercolor{footline}{fg=black!60}
\setbeamertemplate{footline}{%
  \leavevmode\hbox{%
    \begin{beamercolorbox}[wd=.85\paperwidth,ht=2.5ex,dp=1ex,left,leftskip=1em]{footline}\footnotesize\footerlabel\end{beamercolorbox}%
    \begin{beamercolorbox}[wd=.15\paperwidth,ht=2.5ex,dp=1ex,right,rightskip=1em]{footline}\footnotesize\insertframenumber\end{beamercolorbox}}%
  \vskip0pt}
\setbeamertemplate{itemize items}[circle]

\title{Algorithms}
\subtitle{Lecture 1: Growth of Functions and Big-O}
\author{Demo Lecturer}
\date{Semester 2, 2026}

\begin{document}

\begin{frame}
  \titlepage
\end{frame}

\begin{frame}{Housekeeping}
  \begin{itemize}
    \item Weekly quiz: released after the Friday lecture, due the following Friday
    \item Tutorial: Wednesday 15:00 -- bring your attempts at the tute questions
    \item Discussion board and consultation hours: see the course page
    \item Textbook: any standard algorithms textbook (readings are listed per lecture)
  \end{itemize}
\end{frame}

\begin{frame}{Why Do We Care About Efficiency?}
  Two ways to compute $1 + 2 + \dots + n$:
  \medskip
  \begin{columns}[T]
    \column{0.48\textwidth}
    \textbf{Loop}\\[2pt]
    add $1, 2, \dots, n$ one by one\\
    about $n$ additions
    \column{0.48\textwidth}
    \textbf{Formula}\\[2pt]
    compute $n(n+1)/2$\\
    one multiplication, one division
  \end{columns}
  \bigskip
  Both are correct. For $n = 10^9$ only one of them is practical.\\[4pt]
  \alert{Question:} how do we compare algorithms \emph{without} running them?
\end{frame}

\begin{frame}{Input Size and Basic Operation}
  We measure cost as a function of the \textbf{input size} $n$, counting the
  \textbf{basic operation} -- the operation executed most often in the innermost loop.
  \medskip
  \begin{center}
  \begin{tabular}{lll}
    \hline
    Problem & Input size & Basic operation \\
    \hline
    Search a list & number of items $n$ & key comparison \\
    Multiply two matrices & matrix dimension $n$ & scalar multiplication \\
    Test if $n$ is prime & number of bits of $n$ & division \\
    \hline
  \end{tabular}
  \end{center}
\end{frame}

\begin{frame}[fragile]{Worst, Best and Average Case}
\begin{verbatim}
function find(A, x, n)
    j <- 0
    while j < n
        if A[j] = x
            return j
        j <- j + 1
    return -1
\end{verbatim}
  \begin{itemize}
    \item Basic operation: the comparison \verb|A[j] = x|
    \item \textbf{Worst case}: $x$ is absent, so $n$ comparisons
    \item \textbf{Best case}: $x$ is in the first position, so 1 comparison
  \end{itemize}
\end{frame}

\begin{frame}{Big-O: The Definition}
  $O(g(n))$ is the \emph{set} of functions that grow no faster than $g$.
  \begin{block}{Definition}
    $t(n) \in O(g(n))$ iff there exist constants $c > 0$ and $n_0$ such that
    $t(n) \le c \cdot g(n)$ for all $n > n_0$.
  \end{block}
  \textbf{Example.} $3n + 10 \in O(n)$.\\
  For $n > 10$ we have $3n + 10 < 3n + n = 4n$, so choose $c = 4$ and $n_0 = 10$.
\end{frame}

\begin{frame}{Big-Omega and Big-Theta}
  \begin{itemize}
    \item $t(n) \in \Omega(g(n))$: $t(n) \ge c \cdot g(n)$ for all $n > n_0$ -- a \textbf{lower} bound
    \item $t(n) \in \Theta(g(n))$: both $O(g)$ and $\Omega(g)$ -- the \textbf{exact} order of growth
  \end{itemize}
  \bigskip
  \textbf{Caution.} Writing $3n \in O(n^2)$ is true but loose.
  Use $\Theta$ when you can show both bounds.
\end{frame}

\begin{frame}{Common Growth Rates and Summary}
  \[
    1 \;\prec\; \log n \;\prec\; n \;\prec\; n\log n \;\prec\; n^2 \;\prec\; n^3 \;\prec\; 2^n \;\prec\; n!
  \]
  \begin{itemize}
    \item Count the basic operation as a function of input size
    \item Ignore constant factors and small inputs
    \item Next lecture: binary search -- an algorithm with $\Theta(\log n)$ worst case
  \end{itemize}
\end{frame}

\end{document}
